% ============================================================================
% sec3_structural_model.tex — Paper 2: Structural model
% ============================================================================

\section{A Structural Model of Cover-Bidder Deployment}
\label{sec:structural}

We develop a model in which a cartel chooses the number of cover
bidders, their bid distribution, and which markets to target. The
model has three aims: (i)~to formalize the economic logic behind
the frequent-loser phenomenon; (ii)~to specify a mixture likelihood
that separates cover bids from genuine bids in the data; and
(iii)~to derive the \emph{dispersion paradox}---that coordinated
cover bidding produces lower, not higher, within-tender bid
dispersion---as an equilibrium prediction rather than an empirical
curiosity.

% ── 3.1 Environment ──────────────────────────────────────────────

\subsection{Environment}
\label{sec:environment}

Consider a procurement auction for item $i$ at purchasing unit $k$ in
year $t$. The procuring agency posts a reference price $r > 0$ and
selects the lowest bidder.%
\footnote{Under BEC's procurement rules, the reference price is
set by the purchasing unit based on market research---typically
the median of three price quotations from suppliers or the
average of prices in recent comparable tenders. The reference
price is public information, posted before bidding opens.}
There are $n \geq 1$ \emph{genuine} bidders,
each with private cost $c_j \sim F(\cdot \mid \mathbf{x}_j)$, where
$\mathbf{x}_j$ is a vector of observable cost shifters (firm size, CNAE
sector, geographic distance to PBU~$k$). Genuine bidders submit bids
$b_j = c_j + \mu_j$, where $\mu_j \geq 0$ is a markup.

A cartel controls firm $j^*$---the designated winner---and deploys
$m \geq 0$ \emph{cover bidders} (frequent losers). Each cover bidder
submits a bid $b_\ell$ drawn from a distribution $G(\cdot \mid b^*,
\boldsymbol{\theta})$, where $b^*$ is the cartel's designated winning
bid and $\boldsymbol{\theta}$ is a vector of strategic parameters. Cover
bids satisfy $b_\ell > b^*$ with probability
one. The procuring agency observes $n + m$ total bids.

Two institutional features of S\~ao Paulo's BEC platform shape the
model:

\begin{enumerate}
  \item \textbf{Convite} (sealed-bid): requires at least three bidders
    (Lei~8.666/93, Art.~22, \S7). If $n < 3$, the cartel must deploy
    $m \geq 3 - n$ cover bidders to satisfy the minimum-bidder
    constraint.
  \item \textbf{Preg\~ao} (electronic reverse auction): no
    minimum-bidder requirement. Bidders submit descending bids in real
    time, allowing cover bidders to observe auction dynamics before
    calibrating their bids.
\end{enumerate}

\noindent We denote the auction modality by $\tau \in \{\text{convite},
\text{preg\~ao}\}$, the minimum-bidder constraint by $\underline{n}(\tau)$
(equal to 3 for convite, 0 for preg\~ao), and the detection probability
by $\theta_k \in [0,1]$, which varies across PBUs and is increasing in
oversight capacity.

% ── 3.2 Cartel's Problem ─────────────────────────────────────────

\subsection{The Cartel's Optimization Problem}
\label{sec:cartel_problem}

The cartel chooses $m$, $b^*$, and $G$ to maximize expected profit
net of cover-bidding costs and detection penalties:
\begin{equation}
  \max_{m \in \mathbb{Z}_+,\, b^*,\, G} \; \underbrace{\pi(b^*, m, n)
  }_{\text{cartel surplus}}
  - \underbrace{C(m)}_{\text{cover cost}}
  - \underbrace{\theta_k \cdot \Phi(m)}_{\text{expected penalty}}
  \label{eq:cartel_objective}
\end{equation}
subject to:
\begin{align}
  b^* &\leq r \label{eq:ref_price_constraint} \\
  \Pr(b_\ell > b^* \mid G) &= 1 \quad \forall \, \ell \in \{1, \ldots, m\}
    \label{eq:cover_above} \\
  n + m &\geq \underline{n}(\tau) \label{eq:min_bidder}
\end{align}

\noindent We maintain the following assumptions:

\begin{description}
  \item[A1 (Surplus structure).] $\pi(m, n) = b^*(m, n) - c_{j^*}$
    where the sustainable winning bid $b^*$ is increasing in $m$
    ($\partial b^* / \partial m \geq 0$): additional cover bidders
    help sustain a higher $b^*$ by creating the appearance of
    competition.

  \item[A2 (Diminishing returns).] $\partial^2 \pi / \partial m^2 < 0$:
    the marginal return to cover bidding is decreasing.

  \item[A3 (Cover-bidding cost).] $C(m) = c_0 + c_1 m$ with $c_0 \geq 0$
    and $c_1 > 0$: linear in the number of cover bidders. The
    per-bidder cost $c_1 = \kappa + \bar{u}_{(m)}$ is microfounded
    in Section~\ref{sec:cover_participation} as the sum of the
    institutional participation cost and the marginal cover bidder's
    opportunity cost.

  \item[A4 (Detection penalty).] $\theta_k \cdot \Phi(m)$ with
    $\Phi'(m) > 0$ and $\Phi''(m) \geq 0$: more cover bidders increase
    the forensic footprint. For closed-form results we use
    $\Phi(m) = \phi_0 m$ with $\phi_0 > 0$.

  \item[A5 (Interior solution exists).] $\partial \pi / \partial m
    \big|_{m=0} > c_1 + \theta_k \phi_0$.
\end{description}

\begin{proposition}[Optimal number of cover bidders]
\label{prop:optimal_m}
Under Assumptions A1--A5, the cartel's optimal cover-bidder count is:
\begin{equation}
  m^* = \max\!\left\{\underline{n}(\tau) - n, \;
  \left\lfloor \frac{\partial \pi / \partial m}{c_1 + \theta_k \phi_0}
  \right\rfloor \right\}
  \label{eq:optimal_m}
\end{equation}
The optimal $m^*$ satisfies:
\emph{(i)}~$\partial m^* / \partial \theta_k < 0$\,;
\emph{(ii)}~$\partial m^* / \partial c_1 < 0$\,;
\emph{(iii)}~$m^*$ is weakly larger under convite than preg\~ao,
holding $n$ fixed.
\end{proposition}

\noindent\emph{Proof.} See Appendix~\ref{app:proofs}. \hfill$\square$

\medskip

Property~(iii) has a direct institutional implication: the
minimum-bidder rule converts an \emph{interior} solution (cartel
chooses $m$ optimally) into a \emph{corner} solution (cartel is
forced to deploy $m \geq 3 - n$). When $n_{\text{genuine}} = 0$,
the cartel must field at least three cover bidders---a pure
cover-bidding tender where no genuine competition exists.

\paragraph{Quantity vs.\ effectiveness.}
Note that (iii) is a prediction about the \emph{number} of cover
bidders, not about their \emph{effectiveness}. The price impact per
cover bidder depends on the coordination regime
(Section~\ref{sec:cover_distributions}): under Regime~2
(preg\~ao), each cover bidder is more effective because she
calibrates her bid with real-time information. The model therefore
permits $m^*_{\text{convite}} > m^*_{\text{preg\~ao}}$ and yet a
larger price association under preg\~ao---a distinction we test in
Section~\ref{sec:results_ols}.

% ── 3.2b Cover-Bidder Participation ────────────────────────────────

\subsection{Cover-Bidder Participation}
\label{sec:cover_participation}

A cover bidder is a firm that submits bids designed to lose. Why
would a rational firm participate in an auction it cannot win? The
cartel must compensate cover bidders to induce participation; we
formalize this as a participation constraint.

Each cover bidder $\ell$ faces an outside option $\bar{u}_\ell \geq 0$
(the value of her next-best use of time) and incurs a per-tender
participation cost $\kappa > 0$ (registration fees, bid preparation,
bid bonds where applicable). The cartel offers a side payment
$s \geq 0$ per tender. Cover bidder $\ell$ participates if and only
if:
\begin{equation}
  s - \kappa \geq \bar{u}_\ell
  \label{eq:participation_constraint}
\end{equation}

The cartel's cost of deploying $m$ cover bidders is therefore:
\begin{equation}
  C(m) = m \cdot s = m \cdot (\kappa + \bar{u}_{(m)})
  \label{eq:cover_cost_micro}
\end{equation}
where $\bar{u}_{(m)}$ is the outside option of the $m$-th cheapest
cover bidder (ordered by $\bar{u}_\ell$). This microfounds
Assumption~A3: the linear cost $c_1 = \kappa + \bar{u}_{(m)}$
reflects both the institutional participation cost and the cover
bidder's opportunity cost.

\begin{prediction}[Cover-bidder rationality]
\label{pred:rationality}
If FL firms are cover bidders compensated by side payments, they
should exhibit: (i)~low opportunity costs (small firms, narrow
product scope, geographically proximate to target PBUs);
(ii)~persistent relationships with the same set of winners
(reflecting cartel membership); and (iii)~no economic rationale
for participation absent side payments (zero win rate, no learning
or market-entry motivation).
\end{prediction}

\noindent In Section~\ref{sec:spec_tests}, we document that FL firms
are disproportionately micro-enterprises (73\%) with narrow CNAE
scope, consistent with low $\bar{u}_\ell$. Their winner-HHI is 2.4
times the non-FL median, indicating concentrated co-bidding
relationships consistent with cartel membership. FL persistence is
low (10\% across years), consistent with cover bidding as a
transient role assigned by the cartel rather than a permanent firm
characteristic.

\paragraph{Alternative interpretations.}
Two non-cartel explanations for FL behavior merit discussion.
\emph{First}, firms may bid repeatedly to learn about cost
structures or gain experience (``bidding to learn''). This
explanation predicts that FL firms should eventually win some
tenders or exit; the zero win rate over 11~years and high
persistence of participation are inconsistent with a learning
motive. \emph{Second}, firms may bid to maintain ``active bidder''
status in procurement registries. BEC does not require active
bidding for registration renewal, ruling out this mechanism in
our setting.

% ── 3.3 Cover-Bid Distributions ──────────────────────────────────

\subsection{Equilibrium Cover-Bid Distributions}
\label{sec:cover_distributions}

The cartel chooses $G$ to minimize detection probability while
satisfying constraint~\eqref{eq:cover_above}. The key primitive that
distinguishes the two regimes is the \emph{information structure}:
what does a cover bidder know about $b^*$ when submitting a bid? We
formalize two polar cases, corresponding to the two institutional
modalities in BEC.

\begin{description}
  \item[I1 (Limited information).] The cover bidder receives the
    instruction ``bid above $b^*$'' but does not observe $b^*$
    directly. She knows only the reference price $r$ and an upper
    bound $\delta$ on how far above $b^*$ is acceptable (to avoid
    suspiciously high bids). This corresponds to sealed-bid
    (convite) tenders, where bids are simultaneous and $b^*$
    cannot be observed before submission.

  \item[I2 (Full information).] The cover bidder observes $b^*$ in
    real time---or receives a precise target from the cartel
    coordinator---and aims to bid $b^* + \epsilon$ with execution
    noise $\eta \sim N(0, \sigma_c^2)$, $\sigma_c > 0$. This
    corresponds to electronic reverse auction (preg\~ao) tenders,
    where the cartel coordinator can observe the auction dynamics
    and relay a target to cover bidders.
\end{description}

\begin{proposition}[Regime~1: Complementary cover bidding]
\label{prop:regime1}
Under I1, if the cover bidder maximizes entropy over $[b^*, b^* +
\delta]$ subject to $b_\ell > b^*$ (i.e., she has no basis for
concentrating her bid at any particular level), the equilibrium
cover-bid distribution is uniform:
\begin{equation}
  G_1(b) = \frac{b - b^*}{\delta}, \quad b \in [b^*, b^* + \delta]
  \label{eq:regime1}
\end{equation}
with $\text{Var}(b_\ell) = \delta^2 / 12$. If $\delta$ is large
relative to genuine-bid dispersion,
$\sigma_{\text{cover}} > \sigma_{\text{genuine}}$.
\end{proposition}

\begin{proposition}[Regime~2: Coordinated cover bidding]
\label{prop:regime2}
Under I2, the cover bidder targets $b^* + \epsilon$ with execution
noise $\eta$. The resulting cover-bid distribution is truncated
normal:
\begin{equation}
  G_2(b) = \frac{\phi\!\left(\frac{b - (b^* + \epsilon)}
  {\sigma_c}\right)}{\sigma_c\left[1 - \Phi\!\left(
  \frac{b^* - (b^* + \epsilon)}{\sigma_c}\right)\right]},
  \quad b > b^*
  \label{eq:regime2}
\end{equation}
with $\mu_c = b^* + \epsilon$ and truncation at $b^*$ (the cover
bid must exceed the winning bid). Because the truncation compresses
the distribution from below and the target precision $\sigma_c$
reflects the coordinator's real-time information, the effective
dispersion satisfies $\text{sd}(b_\ell \mid b_\ell > b^*) <
\sigma_g$ whenever $\sigma_c \leq \sigma_g$---i.e., coordination
reduces dispersion below the genuine-bid level.
\end{proposition}

\noindent\emph{Proofs.} See Appendix~\ref{app:proofs}. \hfill$\square$

\medskip

The information structure generates a sharp, testable prediction
about the direction of bid dispersion. Under I1 (Regime~1), the
cover bidder's ignorance of $b^*$ forces dispersion \emph{above}
the genuine level. Under I2 (Regime~2), the coordinator's
real-time information allows dispersion \emph{below} it. The two
regimes are empirically distinguishable because they generate
opposite predictions for $\sigma_c / \sigma_g$.

\paragraph{Institutional mapping.}
The information structures map directly onto BEC's two modalities.
Convite (sealed-bid) provides limited information to cover bidders,
as all bids are submitted simultaneously; this favors Regime~1.
Preg\~ao (electronic reverse auction) provides full real-time
information, enabling the precise coordination of Regime~2. If
Regime~2 dominates empirically, it implies that the preg\~ao format
---designed to enhance competition---inadvertently facilitates the
most effective form of cover bidding.

% ── 3.4 Market Selection ─────────────────────────────────────────

\subsection{Market Selection}
\label{sec:market_selection}

\begin{proposition}[Market-selection prediction]
\label{prop:market_selection}
Under Assumptions A1--A2, if the marginal return to cover bidding is
decreasing in market concentration:
\begin{equation}
  \frac{\partial^2 \pi}{\partial m \, \partial \text{HHI}} < 0
  \label{eq:cross_deriv}
\end{equation}
then $m^*$ is decreasing in HHI. Cover bidding concentrates in
competitive markets.
\end{proposition}

\noindent\emph{Proof.} See Appendix~\ref{app:proofs}. \hfill$\square$

\medskip

This reverses the standard detection heuristic. Collusion is commonly
associated with concentrated markets; Proposition~\ref{prop:market_selection}
shows that cover bidding---the operational mechanism of collusion in
procurement---concentrates in \emph{competitive} markets where
manufactured competition is most valuable.

\paragraph{Implication for detection.}
Cover bidding is a strategic complement to genuine competition, not
a substitute. Cartels deploy more cover bidders in competitive
markets (low HHI), where genuine bidders exert downward price
pressure that cover bidders help offset. When HHI is high, the
cartel can sustain prices through market power alone, reducing the
marginal return to synthetic competition.

% ── 3.5 Structural Likelihood ────────────────────────────────────

\subsection{Structural Likelihood}
\label{sec:likelihood}

Each bid $b_{jt}$ in tender $t$ is generated by one of two processes:
\begin{equation}
  b_{jt} \sim
  \begin{cases}
    F_{\text{genuine}}(b \mid \mathbf{x}_j, \boldsymbol{\alpha}) &
      \text{if firm } j \text{ is genuine (non-FL)} \\
    G_{\text{cover}}(b \mid b^*_t, \boldsymbol{\gamma}) &
      \text{if firm } j \text{ is FL}
  \end{cases}
  \label{eq:mixture}
\end{equation}
The FL classification provides the mixture labels, so no EM algorithm
is required---the model is a supervised mixture with observed group
membership.

\paragraph{Genuine bids.} Log-normal with item and year fixed effects:
\begin{equation}
  \log b_{jt} \mid j \text{ genuine} \sim
  N\!\left(\mathbf{x}_j'\boldsymbol{\beta} + \alpha_i + \lambda_t,
  \; \sigma_g^2\right)
  \label{eq:genuine_dist}
\end{equation}
where $\mathbf{x}_j$ includes firm size, firm age, and CNAE sector.

\paragraph{Cover bids.} Under Regime~1:
$b_{jt} \mid j \text{ FL} \sim \text{Uniform}[b^*_t, b^*_t + \delta]$.
Under Regime~2:
$\log b_{jt} \mid j \text{ FL} \sim N(\log b^*_t + \epsilon, \sigma_c^2)$
with $\sigma_c < \sigma_g$.

\paragraph{Log-likelihood.}
\begin{equation}
  \ell_t = \sum_{j \in \mathcal{G}_t} \log f_{\text{genuine}}(b_{jt}
  \mid \mathbf{x}_j, \boldsymbol{\alpha})
  + \sum_{\ell \in \mathcal{F}_t} \log g_{\text{cover}}(b_{\ell t}
  \mid b^*_t, \boldsymbol{\gamma})
  \label{eq:loglik}
\end{equation}

\paragraph{Robustness to classification error.}
The FL classification is a statistical rule that necessarily involves
misclassification. In a supervised mixture with noisy labels,
misclassification attenuates the estimated distributional difference
between groups: if a fraction $\alpha$ of FL-labeled bids are
actually genuine, $\hat{\epsilon}_{\text{noisy}} =
(1-\alpha) \epsilon$. Our estimates of $\hat{\epsilon}$ and
$\hat{\sigma}_c$ are therefore conservative (biased toward the
genuine distribution).

\paragraph{Estimation.}
We estimate $(\boldsymbol{\alpha}, \boldsymbol{\gamma})$ by maximum
likelihood using 23.2~million non-FL losing bids (genuine
parameters) and 189,381 FL losing bids (cover parameters,
conditional on $b^*_t$). Standard errors: parametric bootstrap
(500 replications). Model selection between regimes uses BIC and
Kolmogorov--Smirnov goodness-of-fit.

% ── 3.6 Identification ───────────────────────────────────────────

\subsection{Identification}
\label{sec:identification}

\begin{itemize}
  \item $(\boldsymbol{\beta}, \sigma_g^2)$: identified from non-FL
    losing bids within item--year cells. Item FE absorb
    product-specific cost levels; year FE absorb aggregate trends.

  \item $(\delta \text{ or } \sigma_c, \epsilon)$: identified from
    the distribution of FL bids relative to $b^*_t$ in each tender.

  \item $\theta_k$ (detection probability): partially identified from
    the cross-section of PBU characteristics via
    Proposition~\ref{prop:optimal_m}: $m^*$ is decreasing in
    $\theta_k$, proxied by PBU procurement volume.

  \item Price association: \emph{not} independently identified by
    $(\hat{\epsilon}, \hat{\sigma}_c)$. These characterize the
    cover-bid mechanism, not the price-setting decision. The
    FL--price association comes from the reduced-form comparison of
    FL-present and FL-absent tenders; it reflects both treatment
    and selection, disciplined by the model's predictions.
\end{itemize}

% ── 3.7 Testable Predictions ─────────────────────────────────────

\subsection{Testable Predictions}
\label{sec:predictions}

The model generates five predictions, each tested on independent data
or moments not used in estimation:

\begin{prediction}[Price effect]
\label{pred:price}
FL-present tenders have higher winning prices than comparable tenders
without FL, controlling for item, year, and PBU.
\end{prediction}

\begin{prediction}[Strategic selection]
\label{pred:selection}
FL-present tenders have at least as many genuine (non-FL) competitors
as FL-absent tenders: cover bidders complement, not substitute.
\end{prediction}

\begin{prediction}[Regime identification]
\label{pred:regime}
Under Regime~1, $\sigma_c > \sigma_g$. Under Regime~2,
$\sigma_c < \sigma_g$. BIC selects the dominant regime.
\end{prediction}

\begin{prediction}[Market selection]
\label{pred:market}
The FL--price association is decreasing in market concentration (HHI):
cover bidding concentrates in competitive markets.
\end{prediction}

\begin{prediction}[Exchangeability violation]
\label{pred:exchange}
FL bid residuals are non-exchangeable with non-FL residuals:
the structural model implies FL bids are drawn from $G_{\text{cover}}$,
a distinct data-generating process.
\end{prediction}
